\documentclass[a4paper,12pt]{article}

% ------------------------------------------------
% Кодировка, язык, математика
% ------------------------------------------------
\usepackage{iftex}
\ifPDFTeX
  \usepackage[T2A]{fontenc}
  \usepackage[utf8]{inputenc}
  \usepackage[russian]{babel}
\else
  \usepackage{fontspec}
  \usepackage{polyglossia}
  \setdefaultlanguage{russian}
  \defaultfontfeatures{Ligatures=TeX}
  \setmainfont{DejaVu Sans}
  \setsansfont{DejaVu Sans}
  \AtBeginDocument{\raggedright}
\fi
\usepackage{amsmath,amssymb,mathtools}
\usepackage{icomma}

% ------------------------------------------------
% Шрифт и типографика
% ------------------------------------------------
\renewcommand{\familydefault}{\sfdefault}
\usepackage{microtype}
\usepackage{setspace}
\setstretch{1.12}
\usepackage{parskip}
\setlength{\parindent}{0pt}
\emergencystretch=3em

% ------------------------------------------------
% Страница
% ------------------------------------------------
\usepackage[
  a4paper,
  left=20mm,
  right=20mm,
  top=18mm,
  bottom=21mm
]{geometry}

% ------------------------------------------------
% Списки, таблицы, графики
% ------------------------------------------------
\usepackage{enumitem}
\usepackage{tabularx,array}
\usepackage{tikz}
\usepackage{pgfplots}
\pgfplotsset{compat=1.18}

% ------------------------------------------------
% Цвет и заголовки
% ------------------------------------------------
\usepackage{xcolor}
\definecolor{accent}{RGB}{53,72,92}
\definecolor{soft}{RGB}{105,116,127}

\usepackage{titlesec}
\titleformat{\section}
  {\Large\bfseries\color{accent}}
  {}{0pt}{}
\titlespacing*{\section}{0pt}{0.5em}{0.9em}

% ------------------------------------------------
% Колонтитулы
% ------------------------------------------------
\usepackage{fancyhdr}

\pagestyle{fancy}
\fancyhf{}
\fancyhead[L]{\small\color{soft}Входная диагностическая работа}
\fancyhead[R]{\small\color{soft}10 класс}
\fancyfoot[C]{\small\color{soft}\thepage}

\renewcommand{\headrulewidth}{0pt}
\renewcommand{\footrulewidth}{0pt}
\setlength{\headheight}{15pt}

% ------------------------------------------------
% Ссылки без видимого оформления
% ------------------------------------------------
\usepackage[hidelinks]{hyperref}

% ------------------------------------------------
% Пользовательские настройки
% ------------------------------------------------
\setlist[enumerate]{
  label=\textbf{\arabic*.},
  leftmargin=1.1cm,
  itemsep=0.95em,
  topsep=0.4em
}

\newcommand{\taskspace}{\par\vspace{0.55cm}}
\newcommand{\smalltaskspace}{\par\vspace{0.3cm}}

% ================================================================
% ДОКУМЕНТ
% ================================================================

\begin{document}

% ================================================================
% ТИТУЛЬНЫЙ ЛИСТ
% ================================================================

\begin{titlepage}
\thispagestyle{empty}

\vspace*{26mm}

{\color{soft}\large
Математика \textbullet\ 10 класс
\par}

\vspace{18mm}

{\Huge\bfseries\color{accent}
Входная диагностическая работа
\par}

\vspace{8mm}

{\Large
18 сентября 2026 года
\par}

\vfill

{\large
Лёвин Артём Александрович
\par}

\vspace{2mm}

{\large\bfseries
эксклюзивно для Григория Архипова
\par}

\vspace*{18mm}

\end{titlepage}

% ================================================================
% ПОРЯДОК РАБОТЫ
% ================================================================

\thispagestyle{fancy}

\section*{Порядок работы}

Представленная диагностическая работа позволит оценить наличие базовых
знаний и навыков, необходимых для дальнейшей подготовки к профильному ЕГЭ
по математике.

Пожалуйста, решайте задания последовательно и по возможности оформляйте
решение развёрнуто: записывайте основные преобразования, формулы и
рассуждения, а в конце каждого задания отдельно указывайте полученный ответ.

Если какое-либо задание пока решить не удаётся, достаточно поставить
напротив его номера прочерк и перейти дальше. Это тоже полезная информация
для диагностики: возвращаться к такому заданию наугад или подбирать ответ
не требуется.

Во время работы желательно обходиться без подсказок, справочных материалов
и калькулятора. Если задача занимает слишком много времени, можно пропустить
её и вернуться к ней позднее.

\vspace{8mm}

{\large\bfseries\color{accent}
Структура работы}

\begin{itemize}[
  leftmargin=7mm,
  itemsep=0.45em
]
  \item планиметрия и векторы;
  \item теория вероятностей и статистика;
  \item алгебра и прикладные задачи;
  \item задания повышенной сложности.
\end{itemize}

\vfill

{\small\color{soft}
Цель работы --- определить стартовый уровень и подобрать дальнейшую
траекторию обучения.}

\newpage

% ================================================================
% ПЛАНИМЕТРИЯ И ВЕКТОРЫ
% ================================================================

\section*{Блок I. Планиметрия и векторы}

\begin{enumerate}

\item Через концы $A$ и $B$ дуги окружности с центром $O$ проведены
касательные $AC$ и $BC$. Угол $CAB$ равен $32^\circ$.

Найдите угол $AOB$. Ответ дайте в градусах.
\taskspace

\item Сторона $AB$ треугольника $ABC$ с тупым углом $C$ равна радиусу
описанной около него окружности.

Найдите угол $C$. Ответ дайте в градусах.
\taskspace

\item Даны векторы
\[
\vec a=(1;1),
\qquad
\vec b=(0;2).
\]
Используя скалярное произведение, найдите угол между векторами
$\vec a$ и $\vec b$. Ответ дайте в градусах.
\taskspace

\item Даны векторы
\[
\vec a=(1;1),
\qquad
\vec b=(0;7).
\]
Найдите длину вектора $5\vec a+\vec b$.
\taskspace

\end{enumerate}

\newpage

% ================================================================
% ТЕОРИЯ ВЕРОЯТНОСТЕЙ И СТАТИСТИКА
% ================================================================

\section*{Блок II. Теория вероятностей и статистика}

\begin{enumerate}[resume]

\item В сборнике билетов по биологии всего $35$ билетов, в $14$ из них
встречается вопрос по теме «Ботаника».

Найдите вероятность того, что в случайно выбранном на экзамене билете
школьнику достанется вопрос по теме «Ботаника».
\taskspace

\item На олимпиаде по физике $450$ участников разместили в трёх
аудиториях. В первых двух удалось разместить по $180$ человек,
оставшихся перевели в запасную аудиторию в другом корпусе.

Найдите вероятность того, что случайно выбранный участник писал
олимпиаду в запасной аудитории.
\taskspace

\item Стрелок стреляет по одному разу в каждую из четырёх мишеней.
Вероятность попадания в мишень при каждом отдельном выстреле равна
$0{,}9$.

Найдите вероятность того, что стрелок попадёт в первую мишень и
не попадёт в три последние.
\taskspace

\item При подозрении на наличие некоторого заболевания пациента
отправляют на ПЦР-тест. Если заболевание действительно есть, тест
подтверждает его в $86\%$ случаев. Если заболевания нет, тест выявляет
его отсутствие в среднем в $94\%$ случаев. Известно, что в среднем тест
оказывается положительным у $10\%$ пациентов, направленных на
тестирование.

Тест некоторого пациента оказался положительным.

Какова вероятность того, что пациент действительно имеет это заболевание?
\taskspace

\item В таблице показано распределение случайной величины $X$.

\begin{center}
\renewcommand{\arraystretch}{1.35}
\begin{tabular}{|>{\bfseries}c|c|c|c|c|}
\hline
Значения $X$ & $-5$ & $2$ & $4$ & $5$ \\
\hline
Вероятности & $0{,}2$ & $0{,}4$ & $0{,}3$ & $0{,}1$ \\
\hline
\end{tabular}
\end{center}

Найдите $\mathop{\mathrm{E}}X$ --- математическое ожидание этой
случайной величины.
\taskspace

\end{enumerate}

\newpage

% ================================================================
% АЛГЕБРА И ПРИКЛАДНЫЕ ЗАДАЧИ
% ================================================================

\section*{Блок III. Алгебра и прикладные задачи}

\begin{enumerate}[resume]

\item Найдите $61a-11b+50$, если
\[
\frac{2a-7b+5}{7a-2b+5}=9.
\]
\taskspace

\item Найдите значение выражения
\[
\left(\frac{3}{4}+2\frac{3}{8}\right)\cdot 25{,}6.
\]
\taskspace

\item Коэффициент полезного действия двигателя определяется формулой
\[
\eta=\frac{T_1-T_2}{T_1}\cdot 100\%,
\]
где $T_1$ --- температура нагревателя в градусах Кельвина,
$T_2$ --- температура холодильника в градусах Кельвина.

При какой минимальной температуре нагревателя $T_1$ КПД этого двигателя
будет не меньше $15\%$, если температура холодильника $T_2=340$ К?
Ответ выразите в градусах Кельвина.
\taskspace

\item Камнеметательная машина выстреливает камни под некоторым острым
углом к горизонту. Траектория полёта камня описывается формулой
\[
y=ax^2+bx,
\qquad
a=-\frac{1}{110}\ \text{м}^{-1},
\qquad
b=\frac{13}{11},
\]
где $x$ --- смещение камня по горизонтали в метрах, $y$ --- высота
камня над землёй в метрах.

На каком наибольшем расстоянии в метрах от крепостной стены высотой
$19$ м нужно расположить машину, чтобы камни пролетали над стеной
на высоте не менее $1$ метра?
\taskspace

\item Поезд, двигаясь равномерно со скоростью $60$ км/ч, проезжает
мимо лесополосы, длина которой равна $400$ метрам, за $1$ минуту.

Найдите длину поезда в метрах.
\taskspace

\newpage

\item Теплоход, скорость которого в неподвижной воде равна $15$ км/ч,
проходит по течению реки и после стоянки возвращается в исходный пункт.
Скорость течения равна $3$ км/ч, стоянка длится $5$ часов, а в исходный
пункт теплоход возвращается через $25$ часов после отплытия из него.

Сколько километров проходит теплоход за весь рейс?
\taskspace

\item На рисунке изображён график функции
\[
f(x)=\frac{k}{x}+a.
\]
Найдите $f(50)$.

\begin{center}
\begin{tikzpicture}
\begin{axis}[
    width=10.6cm,
    height=6.7cm,
    axis lines=middle,
    xmin=-5.5,
    xmax=5.5,
    ymin=-5.5,
    ymax=4.5,
    xtick={-5,-4,-3,-2,-1,0,1,2,3,4,5},
    ytick={-5,-4,-3,-2,-1,0,1,2,3,4},
    grid=both,
    major grid style={draw=black!14},
    minor tick num=0,
    tick label style={font=\small},
    xlabel={$x$},
    ylabel={$y$},
    xlabel style={at={(ticklabel* cs:1)},anchor=north west},
    ylabel style={at={(ticklabel* cs:1)},anchor=south east},
    samples=180,
    clip=true
]
\addplot[dashed,color=soft,thick,domain=-5.5:5.5] {-3};
\addplot[very thick,color=accent,domain=-5.5:-0.24] {2/x-3};
\addplot[very thick,color=accent,domain=0.24:5.5] {2/x-3};
\addplot[only marks,mark=*,mark size=2.3pt,color=accent]
  coordinates {(-2,-4)};
\end{axis}
\end{tikzpicture}
\end{center}
\smalltaskspace

\item На рисунке изображён график функции
\[
f(x)=ax^2+bx+c,
\]
где числа $a$, $b$ и $c$ --- целые. Найдите $f(2)$.

\begin{center}
\begin{tikzpicture}
\begin{axis}[
    width=10.6cm,
    height=6.7cm,
    axis lines=middle,
    xmin=-6.5,
    xmax=1.5,
    ymin=-5,
    ymax=5.5,
    xtick={-6,-5,-4,-3,-2,-1,0,1},
    ytick={-4,-3,-2,-1,0,1,2,3,4,5},
    grid=both,
    major grid style={draw=black!14},
    minor tick num=0,
    tick label style={font=\small},
    xlabel={$x$},
    ylabel={$y$},
    xlabel style={at={(ticklabel* cs:1)},anchor=north west},
    ylabel style={at={(ticklabel* cs:1)},anchor=south east},
    samples=180,
    domain=-6.5:0.5,
    clip=true
]
\addplot[very thick,color=accent] {(x+3)^2-4};
\end{axis}
\end{tikzpicture}
\end{center}
\smalltaskspace

\newpage

\item
\begin{enumerate}[
  label=\alph*),
  leftmargin=8mm,
  itemsep=0.55em,
  topsep=0pt
]
  \item Решите уравнение
  \[
  \frac{(x-1)^2}{8}+\frac{8}{(x-1)^2}
  =
  7\left(\frac{x-1}{4}-\frac{2}{x-1}\right)-1.
  \]
  \item Найдите его корни, принадлежащие отрезку $[-2;3]$.
\end{enumerate}
\taskspace

\item Решите неравенство
\[
\frac{3x+28-x^2}{x^2-7x}
\le
2+\frac{2}{x+3}.
\]
\taskspace

\end{enumerate}

\newpage

% ================================================================
% ЗАДАНИЯ ПОВЫШЕННОЙ СЛОЖНОСТИ
% ================================================================

\section*{Блок IV. Задания повышенной сложности}

{\color{soft}
В заданиях этого блока особенно важно записать полное решение:
основные преобразования, используемые свойства и необходимые обоснования.
}

\vspace{4mm}

\begin{enumerate}[resume]

\item Две окружности касаются внутренним образом в точке $C$.
Вершины $A$ и $B$ равнобедренного прямоугольного треугольника $ABC$
с прямым углом $C$ лежат на большей и меньшей окружностях соответственно.
Прямая $AC$ вторично пересекает меньшую окружность в точке $D$.
Прямая $BC$ вторично пересекает большую окружность в точке $E$.

\begin{enumerate}[
  label=\alph*),
  leftmargin=8mm,
  itemsep=0.35em,
  topsep=0.35em
]
  \item Докажите, что $AE\parallel BD$.
  \item Найдите $AC$, если радиусы окружностей равны $8$ и $15$.
\end{enumerate}
\vspace{1.8cm}

\item Найдите все значения $a$, при каждом из которых уравнение
\[
\sqrt{4^x-a}+\frac{a-3}{\sqrt{4^x-a}}=1
\]
имеет ровно два различных корня.
\vspace{2.0cm}

\item На доске записано $10$ натуральных чисел, среди которых нет
одинаковых. Оказалось, что среднее арифметическое любых трёх, четырёх,
пяти или шести чисел из записанных является целым числом. Одно из
записанных чисел равно $30\,021$.

\begin{enumerate}[
  label=\alph*),
  leftmargin=8mm,
  itemsep=0.35em,
  topsep=0.35em
]
  \item Может ли среди записанных на доске чисел быть число $351$?
  \item Может ли отношение двух записанных на доске чисел равняться $11$?
  \item Отношение двух записанных на доске чисел является целым числом
  $n$. Найдите наименьшее возможное значение $n$.
\end{enumerate}

\vfill

{\small\color{soft}
Работа завершена. Если некоторые задания остались без решения,
оставьте их с прочерком --- это поможет точнее определить темы
для повторения.}

\end{enumerate}

\end{document}
